\section{Theoretisch kader}
\subsection{Master Data Management}
Master Data Management (MDM) wordt hier operationeel afgebakend als de gezamenlijke inrichting van processen, verantwoordelijkheden en voorzieningen voor gegevens over herhaald gebruikte domeinentiteiten. De organisatorische en technische dimensies worden ondersteund door het referentiemodel van Otto et al.; de onderzochte context is bedrijfsgericht en vormt geen directe validatie voor Nederlandse basisregistraties. \citep{L05} Haug et al. onderzochten barrières in 787 Deense productiebedrijven. De abstractbevinding dat organisatorische problemen expliciet aandacht verdienen ondersteunt het opnemen van governance, maar niet een kwantitatieve effectverwachting voor het Stelsel. \citep{L17}

DAMA-DMBOK wordt gebruikt als positionering van datamanagement als discipline. Omdat alleen officiële editie-informatie is vastgelegd, worden geen inhoudelijke hoofdstukclaims of specifieke architectuurvoorschriften aan het boek toegeschreven. \citep{S01} ISO 8000 wordt evenmin als één allesomvattend MDM-model behandeld. De delen 100 en 110 betreffen onder meer de uitwisseling van karakteristieke masterdata en semantic encoding; delen 120, 130 en 140 behandelen respectievelijk provenance, nauwkeurigheid en volledigheid binnen hun scope. \citep{S02,S03,S04,S05,S06} Uit deze toepassingsgebieden volgt niet dat RDF verplicht is, dat interne MDM-processen volledig zijn beschreven of dat een concrete kwaliteitsdrempel is gegeven.

Een \emph{masterdata-entiteit} is in dit ontwerp de referent waarover meerdere processen gegevens gebruiken. Een \emph{masterdata-record} is een geregistreerde representatie van uitspraken over die referent in een bepaald systeem en een bepaalde versie. Een entiteit kan daarmee meerdere records hebben; deze modelkeuze impliceert geen administratieve samenvoegbevoegdheid. \emph{Referentiedata} duiden hier beheerde waardenstelsels aan waarmee andere waarden worden geclassificeerd. Het onderscheid moet per domein worden ingevuld en mag niet uit de technische tabelnaam worden afgeleid.


\subsubsection{Operationele afbakening en grensgevallen}
Binnen dit onderzoek is een publicatie precies dan een SMDP-kandidaat als zij gezamenlijk aan vier voorwaarden voldoet: (a) zij levert of ontsluit herhaald gebruikte uitspraken over geïdentificeerde domeinentiteiten; (b) een producteigenaar legt gebruiksdoel, versie en verantwoordelijkheden vast; (c) geselecteerde gegevenselementen zijn via getypeerde verwijzingen verbonden met hun interpretatiecontext; en (d) het contract beschrijft hoe afnemers bron-, tijd-, gezags- en ontbrekende-bewijsinformatie kunnen verkrijgen. Dit zijn noodzakelijke en gezamenlijk voldoende voorwaarden voor de \emph{onderzoekscategorie}, niet voor kwaliteit, normconformiteit of bewezen nut. Een concrete toepassing kan eraan voldoen en toch de centrale these weerleggen.

Een begrippenregister zonder aangeboden mastergegevens voldoet niet aan (a). Een dataset met alleen een technische schema-URL voldoet niet vanzelf aan (c). Een volledig gedocumenteerd conventioneel register kan wel voldoen: het begrip schrijft geen graafdatabase voor. Het experiment moet daarom de mate en vorm van contextbinding manipuleren in plaats van simpelweg producten met en zonder het label SMDP te vergelijken.

Masterdata, referentiedata en metadata zijn gebruiksrollen. Een beheerde waardenlijst classificeert in de referentierol, terwijl de beschrijving van die lijst metadata is. De organisatie die een product bezit kan zelf onderwerp van masterdata zijn; haar vermelding als eigenaar is productmetadata. Een gebeurtenisrecord is niet louter door hergebruik masterdata. Dataset, dienst en product beschrijven publicatie-/gebruiksobjecten en behoren niet tot dezelfde disjuncte classificatie als inhoudsrollen. Deze operationalisering sluit aan op de onderscheiden registratieniveaus en productoriëntatie, maar is eigen onderzoeksafbakening. \citep{S08,S47,L12,L19}

\subsection{Metadata en data-elementen}
ISO/IEC 11179 onderscheidt registratievoorzieningen van specifieke geregistreerde inhoud. Deel 3 beschrijft algemene voorzieningen; deel 31 richt zich op data-elementen; deel 33 op datasets; deel 6 op registratie. \citep{S07,S08,S09,S10} Voor de paper is vooral van belang dat een data-elementconcept niet samenvalt met een concrete representatie: objectklasse en eigenschap enerzijds en conceptueel domein en waardedomein anderzijds maken verschillende vragen expliciet. \citep{S08}

Een \emph{data-element} wordt hier begrepen als een geregistreerde eenheid waarin de betekenis van een eigenschap wordt verbonden met haar toegestane representatie. Een \emph{attribuut} is een eigenschapconstructie binnen een model; een \emph{waarde} is een concrete invulling in een record. De voorgestelde koppeling van een MIM-attribuut aan een 11179-data-element is daarom een mapping tussen metamodelconstructies, geen vanzelfsprekende identiteit. ISO/IEC 11179-31 betreft bovendien niet alleen data-elementen maar ook bijbehorende concepten, domeinen, eenheden en afleidingsregels; het ontwerp introduceert deze niet als nieuwe constructies. MIM en 11179 bieden aangrenzende aangrijpingspunten, maar het corpus levert geen universele verliesvrije transformatie tussen beide. \citep{S13,S08}

Metadata zijn niet automatisch minder gezaghebbend of minder kritisch dan de beschreven gegevens. In het ontwerp krijgen daarom ook definitie-, model- en mappingversies identifiers, provenance en lifecyclemetadata. Dit is een ontwerpafleiding uit metadata\-registratie en provenance, niet een afzonderlijke eis die hier uit één clausule wordt overgenomen. \citep{S07,S10,S33}

\subsection{Terminologie, definities en concepten}
ISO 704 behandelt de verhouding tussen objecten, begrippen, definities en aanduidingen. NL-SBB biedt een Nederlands model voor het beschrijven van begrippen; SKOS biedt een RDF-vocabulaire voor kennisorganisatiesystemen. \citep{S11,S14,S31} De paper hanteert de werkdefinities in tabel~\ref{tab:concepts}. Zij zijn een analytische afbakening voor het artefact en worden niet gepresenteerd als letterlijke normcitaten.

\begin{longtable}{@{}P{.23\linewidth}P{.70\linewidth}@{}}
\caption{Kernbegrippen en expliciete grenzen. Grondslag: terminologie, metadata\-registratie en semantische modellering. \citep{S11,S08,S13,S28,S30,S31,S47,L12}}\label{tab:concepts}\\
\toprule Begrip & Werkdefinitie en grens \\ \midrule\endfirsthead
\toprule Begrip & Werkdefinitie en grens \\ \midrule\endhead
Term & Talige aanduiding van een begrip in een taal en context; geen identifier op grond van de schrijfwijze alleen. \\
Begrip/concept & Betekeniseenheid waarmee objecten, kenmerken of situaties worden begrepen; in het metamodel onderscheiden van woorden en gegevensrecords. \\
Definitie & Uitspraak die een begrip in een context afbakent; niet hetzelfde als een veldbeschrijving of een opsomming van mogelijke waarden. \\
Juridisch begrip & Begrip waarvan de betekenis of toepassing aan een juridische bron en context wordt verbonden; niet noodzakelijk volledig formaliseerbaar. \\
Conceptueel object & Individuele referent in het gemodelleerde domein; de representatie daarvan is niet het object zelf. \\
Ontologische klasse & Formele categorie in een ontologie; een OWL-klasse specificeert via axioma’s mogelijke leden en relaties, geen registratieplicht. \\
MIM-objecttype & Constructie in een informatiemodel voor objecten met overeenkomstige kenmerken; niet automatisch identiek aan een OWL-klasse of SKOS-concept. \\
Attribuut & Eigenschapconstructie van een gemodelleerd type; geen concrete attribuutwaarde. \\
Data-element & Geregistreerde verbinding van een data-elementconcept met een waarderepresentatie; geen recordcel. \\
Identifier & Aanduiding gebruikt om binnen vastgestelde scope een referent te onderscheiden; scope en uitgiftebeleid zijn onderdeel van de interpretatie. \\
Masterdata-entiteit & Domeinreferent waarvoor herhaald gebruikte kerngegevens worden beheerd; kan administratief of fysiek van aard zijn. \\
Masterdata-record & Geregistreerde verzameling uitspraken over een entiteit; heeft eigen bron, versie en tijdscontext. \\
Dataset & Afgebakende verzameling gegevens als catalogiseerbare eenheid; niet gelijkgesteld aan de specifieke RDF-betekenis van ``dataset''. \\
Data service & Dienst die toegang of verwerking beschikbaar stelt; een endpoint is een toegangspunt van een dienst. \\
Data product & Beheerde eenheid met doel, eigenaar, contract en gekoppelde datasets/diensten; kan meer omvatten dan één dataset. \\
\bottomrule
\end{longtable}

\subsubsection{Definitie versus datastructuur}
Een schema kan aangeven dat \texttt{oppervlakte} een numeriek veld is, eventueel met een eenheid en toegestane grenzen. Daarmee is nog niet noodzakelijk bepaald welk object wordt gemeten, welke afbakeningsregel geldt of op welk tijdstip de waarde betrekking heeft. Dit is een analytisch tegenvoorbeeld: meerdere betekenissen kunnen dezelfde technische vorm hebben. ISO 704 adresseert begripsafbakening, ISO/IEC 11179 de koppeling aan concept en representatie en OpenAPI de beschrijving van een interface. \citep{S11,S08,S19} Hun verschillende scopes ondersteunen de conclusie dat structuur alleen geen garantie voor volledige betekenis biedt.

De omgekeerde reductie is eveneens ongewenst: een natuurlijke-taaldefinitie levert niet vanzelf een uitvoerbare validator. Het ontwerp stelt daarom een expliciete verbinding voor tussen definitie, modelelement en constraint, met documentatie van wat niet is geformaliseerd. De betekenis wordt niet volledig gelijkgesteld aan de verzameling technisch afdwingbare regels.

\subsection{Semantiek en representatie}
RDF definieert een graafgebaseerd gegevensmodel met uitspraken over geïdentificeerde resources. RDFS voegt onder meer klassen en eigenschappen toe; OWL ondersteunt verdergaande ontologische axioma’s. SKOS modelleert begrippenstelsels met labels en semantische relaties. \citep{S28,S29,S30,S31} Deze instrumenten vullen elkaar aan, maar zijn niet onderling verwisselbaar. SKOS §3.5.1 laat dezelfde resource als SKOS-concept en OWL-klasse toe. De afzonderlijke beheerrollen in dit ontwerp zijn dus geen universele disjunctheidsclaim. Meervoudige typering is toegestaan als interpretatie en beheer expliciet blijven. \citep{S31} Een SKOS-label is geen definitie op grond van zijn aanwezigheid; een bredere term is niet automatisch een ontologische superclass; conceptuele matching is geen onbeperkte identiteit van alle bedoelde objecten. \citep{S31,S30}

JSON-LD verbindt JSON-representaties met gekoppelde gegevens via onder meer contexten. SPARQL biedt query- en protocolvoorzieningen voor RDF. Geen van beide bepaalt zelfstandig de domeinbetekenis, het gezag van een bron of de juistheid van een mapping. \citep{S35,S36} De architectuur kiest RDF daarom als kandidaat voor de verbindingsgraaf en JSON-LD als mogelijke distributievorm, zonder deze keuze tot enig juiste opslagarchitectuur te verheffen. Een relationele of documentgebaseerde bron kan blijven bestaan wanneer de semantische verbindingen expliciet en toetsbaar worden gepubliceerd.

\subsubsection{Semantic false agreement}
Deze paper gebruikt \emph{semantic false agreement} als werkterm voor onterecht aangenomen betekenisequivalentie (P05). Er is in het corpus geen afzonderlijk geverifieerde theorie met deze naam. De aanleiding volgt wel uit het onderscheid tussen aanduiding en begrip en uit de verschillende soorten semantische relaties in SKOS en OWL. \citep{S11,S31,S30}

Er worden twee hypothetische fouttypen onderscheiden. Bij schijnovereenstemming delen systemen een label, maar verschillen de referenten, geldigheidsvoorwaarden of definities. Bij gemiste overeenstemming verschillen labels terwijl een voldoende onderbouwde conceptuele overeenkomst bestaat. Het ontwerp moet beide gevallen kunnen onderzoeken zonder automatisch gelijkheid te concluderen. Een mapping krijgt daarom richting, relatie, context, bron- en doelversie en een besluitgrond. Bij onvoldoende bewijs blijft de relatie onzeker of ontbreekt zij; \texttt{owl:sameAs} is geen standaardantwoord op vergelijkbare labels.

\subsection{Ontologie en ontologische commitment}
Een ontologische commitment betreft hier de aannames die het ontwerp doet over soorten entiteiten, hun identiteit en hun afhankelijkheden. Het opnemen van een objecttype suggereert een categorie waarvan instanties kunnen worden onderscheiden; het modelleren van een rol vraagt bovendien onder welke voorwaarden die rol geldt. UFO en ontologiegedreven conceptual modelling bieden een grondslag om dergelijke aannames expliciet te analyseren. OntoClean maakt onder meer identiteit en rigiditeit tot bespreekbare metaeigenschappen. \citep{L08,L09,L10}

Een conceptueel informatiemodel kan zelf semantisch zijn; semantiek is geen exclusief tussenstation vóór ontologie. MIM §1.6 onderscheidt beschouwingsniveaus en formuleert begrip als combinatie van term en definitie. Dit onderzoek beheert de aanduiding, definitie-uitspraak en het bedoelde concept afzonderlijk. Dat is een expliciete operationalisering naast het MIM-woordgebruik, geen stilzwijgend aangenomen equivalentie. \citep{S13}

Een foundational ontology levert algemene onderscheidingen; een domeinontologie specificeert categorieën en relaties binnen een toepassingsgebied. OntoUML is een ontologisch gefundeerde modelleertaal, geen Nederlandse overheidsstandaard. De communitydocumentatie en academische bronnen ondersteunen die positionering, maar leveren in dit dossier geen aangetoonde meerwaarde voor basisregistraties. \citep{S40,S41,L08,L09} MIM is daarentegen een metamodel voor informatiemodellering. Een UFO-analyse kan een MIM-model informeren, maar een MIM-model wordt niet door zijn notatie alleen een foundational ontology. \citep{S13}

De ontwerpkeuze is selectief: ontologische verdieping wordt ingezet waar concurrerende interpretaties van identiteit, rol, fase of afhankelijkheid concrete risico’s veroorzaken. De hypothese G03 is dat deze analyse fouten kan blootleggen die een beperktere modelreview mist. Een vergelijking met een MIM-analyse zonder UFO moet vaststellen of die winst opweegt tegen extra inspanning. Het corpus rechtvaardigt geen universele verplichting om elk begrip in OntoUML te modelleren.

\subsection{Semantische interoperabiliteit}
Het EIF onderscheidt vier interoperabiliteitslagen. Voor de technische analyse in deze paper wordt de technische laag nader onderscheiden in syntactische en structurele interoperabiliteit; dit is een eigen analytische verfijning, geen vijfde officiële EIF-laag. \citep{S43} \emph{Syntactische interoperabiliteit} betekent hier dat representaties parseerbaar zijn. \emph{Structurele interoperabiliteit} betreft overeenkomst over velden, typen en relaties. \emph{Semantische interoperabiliteit} betreft voldoende gedeelde interpretatie van de bedoelde betekenis binnen de gebruikscontext. \emph{Organisatorische interoperabiliteit} betreft afstemming van processen en verantwoordelijkheden; \emph{juridische interoperabiliteit} betreft de verenigbaarheid en toepasselijkheid van juridische kaders.

Een succesvolle parse is daarmee een noodzakelijke stap voor veel vormen van geautomatiseerde uitwisseling, maar geen bewijs van de overige vormen. Een mapping tussen schema’s kan structureel volledig zijn en toch een conceptueel verschil verhullen. Evenmin verleent semantische consistentie een bevoegdheid om persoonsgegevens te verwerken. De architectuur houdt juridische toepasselijkheid daarom als afzonderlijke verantwoordelijkheid bij; deze paper voert geen concrete juridische toets van gegevensverwerkingen uit.

\subsection{Data Products en governance}
De gebruikersoriëntatie in dataproductliteratuur ondersteunt het opnemen van doel, eigenaar, kwaliteit en gebruiksvoorwaarden in het ontwerp. \citep{L12} DCAT 3, DCAT-AP en DCAT-AP-NL bieden daarvoor catalogusmetadata voor onder meer datasets, distributies en diensten. Zij worden niet gelijkgesteld aan een compleet SMDP-contract met afhankelijkheden, operationele afspraken en epistemische regels. \citep{S47,S45,S16} Het aanvullende productcontract is een expliciete ontwerpuitbreiding, waarvan de aansluiting op catalogusprofielen moet worden getoetst.

Governance omvat in deze paper besluitrechten en verantwoordelijkheden voor definities, mappings, kwaliteit en publicatie. Abraham et al. ordenen governance in een conceptueel raamwerk; BOMOS ondersteunt het beheer van open standaarden. \citep{L07,S23} Deze bronnen motiveren beheer als onderdeel van het artefact, maar BOMOS is geen operationeel MDM-schema. Er wordt daarom een onderscheid voorgesteld tussen het beheren van de standaard, het beheren van de toepassing ervan en het beheren van individuele gegevens.

\subsection{Identiteit en persistente identificatie}
URI- en IRI-specificaties geven syntactische regels voor identificatie, terwijl Cool URIs richtlijnen biedt voor het onderscheid tussen resources en hun webdocumenten. \citep{S38,S39} Persistentie is daarmee niet louter een geldige tekenreeks: het ontwerp moet uitgifte, scope, resolutie, overdracht en uitfasering organiseren. Een entiteitidentifier, recordidentifier, versieidentifier en document-URL krijgen afzonderlijke rollen (R03).

Of twee registraties hetzelfde object betreffen, is een inhoudelijke vraag. Gelijke namen of coördinaten volstaan in het voorgestelde model niet als identiteitsbewijs; een geïdentificeerde koppelregel en haar rechts- en modelcontext zijn nodig. Ontologische analyse kan de gehanteerde identiteitscriteria expliciteren. \citep{L08,L10} Het ontwerp ondersteunt daarom zowel vastgestelde identiteit als zwakkere relaties zoals mogelijke correspondentie of betrekking op dezelfde locatie, zonder deze onder één gelijkheidsrelatie samen te voegen.

\subsection{Temporaliteit en historie}
Valid time betreft de tijd waarin een feit in de gemodelleerde werkelijkheid geldig is; transaction time betreft de tijd waarin het feit in de database is vastgelegd. \citep{L14} De term \emph{registratietijd} wordt in het ontwerp alleen als transaction time gebruikt als de betekenis en het betreffende registratiesysteem expliciet zijn vastgelegd. Een ontvangstmoment bij een afnemer is niet automatisch de registratietijd bij de bron.

OWL-Time kan tijdstippen en intervallen representeren, maar schrijft geen complete bitemporele registratiearchitectuur voor. Het tijdveld van een CloudEvent vervangt evenmin een brongebonden geldigheidsinterval. \citep{S37,S22} R06 onderscheidt daarom geldigheid, registratie, publicatie en eventtijd. Bij een latere correctie kan de geldigheid terugwerken terwijl het kennismoment verandert. Oude antwoorden moeten dan kunnen worden gereconstrueerd voor zowel ``geldig op'' als ``bekend bij systeem X op''. Dit is een voorgestelde ontwerpverplichting, geen bewering dat het prototype dit reeds kan.

\subsection{Provenance, authenticiteit en epistemische status}
PROV-O beschrijft relaties tussen entiteiten, activiteiten en agents; ISO 8000-120 adresseert provenance binnen masterdata-uitwisseling. \citep{S33,S04} Provenance maakt een herkomstclaim uitdrukbaar, maar verifieert die claim niet zelfstandig. Een technisch geauthenticeerde afzender levert evenmin vanzelf bewijs van wettelijke authenticiteit of feitelijke juistheid. FSC adresseert vertrouwen en uitwisseling tussen deelnemende systemen; de basisregistratie-eisen betreffen daarnaast institutionele verantwoordelijkheden. \citep{S21,S25}

Het ontwerp onderscheidt vijf door de onderzoeksvraag vereiste categorieën: een normatieve definitie, een authentiek of anderszins gezaghebbend geregistreerd gegeven, een niet-authentiek geregistreerd gegeven, een deterministisch afgeleid gegeven en een probabilistisch of AI-afgeleid gegeven. Deze categorieën mengen echter verschillende assen. Een deterministisch afgeleide waarde kan ook geregistreerd worden; een normatieve definitie kan in een register staan. Daarom stelt R23 geen exclusieve rangorde voor, maar een samengesteld metadatarecord met \emph{soort uitspraak}, \emph{gezagsgrond}, \emph{afleidingswijze}, \emph{bewijsstatus} en \emph{verantwoordelijke}.

Wettelijke authenticiteit wordt apart vastgelegd van bestuurlijk gezag. Een onbekende status wordt niet als ``niet-authentiek'' gecodeerd. Een probabilistische score vormt geen rechtsgrond; een deterministische berekening kan onjuiste invoer gebruiken. Dit zijn ontwerpregels om onterechte gevolgtrekkingen te blokkeren. De bronnen ondersteunen de bouwstenen herkomst, gezag en machinegericht hergebruik, maar niet de volledige samengestelde classificatie. \citep{S25,S33,L11,L16} Of deze metadata in de praktijk betrouwbaar kan worden gevuld en correct wordt gebruikt, wordt in vervolgonderzoek getoetst.

\subsection{Kwaliteit, constraints en AI als machineconsument}
Datakwaliteit omvat meer dan nauwkeurigheid. ISO/IEC 25012 biedt kwaliteitskenmerken, Wang en Strong onderzoeken het perspectief van datagebruikers en NORA biedt een overheidsgericht kwaliteitsraamwerk. \citep{S12,L15,S48} DQV representeert kwaliteitsmetadata; SHACL specificeert constraints op een gegevensgraaf. \citep{S34,S32} Een SHACL-conform rapport betekent daarom uitsluitend dat de onderzochte graaf onder de gekozen validatiecondities aan de toegepaste shapes voldoet. Het bewijst geen feitelijke juistheid in de werkelijkheid.

Ook inferentie en validatie moeten worden gescheiden. RDFS-domein- en bereikuitspraken ondersteunen gevolgtrekkingen; OWL hanteert logische semantiek die niet gelijkstaat aan de gesloten invulcontrole van een formulier. \citep{S29,S30} De validator moet dus vermelden welke gegevens, shapeversie en eventuele inferentiestap zijn gebruikt. De kwaliteitsmeting moet daarnaast populatie, tijdvenster, methode en maatstaf aangeven. De combinatie is een ontwerpafleiding uit de verschillende scopes.

FAIR biedt theoretische steun voor machinevindbare en herbruikbare data en metadata. Kennisgraafliteratuur beschrijft expliciete modellen en verwerkingstechnieken, terwijl Lewis et al. voor specifieke NLP-taken voordelen van retrieval-augmented generation rapporteren. \citep{L11,L13,L16} Deze bronnen onderbouwen de mogelijkheid om relevante context machinegericht beschikbaar te stellen. Zij onderbouwen geen algemene causale wet dat meer ontologie altijd tot minder AI-fouten leidt.

De voorgestelde hypothese is daarom conditioneel: wanneer een consument relevante semantische informatie kan vinden, correct kan verwerken en op de juiste taak kan toepassen, kunnen expliciete betekenis- en bewijsrelaties de interpretatie verbeteren (G02/P02). Verkeerde mappings, ongeschikte contextselectie of een consument die metadata negeert kunnen dat effect neutraliseren. De evaluatie moet vaststellen of het SMDP voordelen biedt boven dezelfde broninhoud zonder expliciete koppelingen, en of eventuele winst samenhangt met semantiek, extra informatie of technische toegang.
